\documentclass[11pt]{article}

% ===================================================================
% Shared preamble for the "tiny lemma, read slowly" preprint series.
% Mirrors the style of FrobeniusLemmas_preprint.tex.
% ===================================================================

\usepackage{amsmath,amssymb,amsthm}
\usepackage{mathtools}
\usepackage[margin=1.2in]{geometry}
\usepackage{hyperref}
\usepackage{xcolor}
\usepackage{listings}
\usepackage{tikz}
\usepackage{tikz-cd}
\usepackage{mathpartir}
\usepackage{newunicodechar}
\usepackage{setspace}

\onehalfspacing

% ---- Unicode glyphs ----
\newunicodechar{→}{\ensuremath{\to}}
\newunicodechar{↦}{\ensuremath{\mapsto}}
\newunicodechar{↔}{\ensuremath{\leftrightarrow}}
\newunicodechar{⟹}{\ensuremath{\Longrightarrow}}
\newunicodechar{⟶}{\ensuremath{\longrightarrow}}
\newunicodechar{≠}{\ensuremath{\neq}}
\newunicodechar{≤}{\ensuremath{\leq}}
\newunicodechar{≥}{\ensuremath{\geq}}
\newunicodechar{∀}{\ensuremath{\forall}}
\newunicodechar{∃}{\ensuremath{\exists}}
\newunicodechar{∘}{\ensuremath{\circ}}
\newunicodechar{·}{\ensuremath{\cdot}}
\newunicodechar{∈}{\ensuremath{\in}}
\newunicodechar{∉}{\ensuremath{\notin}}
\newunicodechar{≃}{\ensuremath{\simeq}}
\newunicodechar{≅}{\ensuremath{\cong}}
\newunicodechar{≈}{\ensuremath{\approx}}
\newunicodechar{∧}{\ensuremath{\wedge}}
\newunicodechar{∨}{\ensuremath{\vee}}
\newunicodechar{¬}{\ensuremath{\neg}}
\newunicodechar{⊢}{\ensuremath{\vdash}}
\newunicodechar{⟨}{\ensuremath{\langle}}
\newunicodechar{⟩}{\ensuremath{\rangle}}
\newunicodechar{⇑}{\ensuremath{\mathord{\uparrow}}}
\newunicodechar{↑}{\ensuremath{\mathord{\uparrow}}}
\newunicodechar{∎}{\ensuremath{\blacksquare}}
\newunicodechar{∣}{\ensuremath{\mid}}
\newunicodechar{∤}{\ensuremath{\nmid}}
\newunicodechar{∑}{\ensuremath{\textstyle\sum}}
\newunicodechar{∏}{\ensuremath{\textstyle\prod}}
\newunicodechar{∅}{\ensuremath{\emptyset}}
\newunicodechar{∪}{\ensuremath{\cup}}
\newunicodechar{∩}{\ensuremath{\cap}}
\newunicodechar{≡}{\ensuremath{\equiv}}
\newunicodechar{ℕ}{\ensuremath{\mathbb{N}}}
\newunicodechar{ℤ}{\ensuremath{\mathbb{Z}}}
\newunicodechar{ℝ}{\ensuremath{\mathbb{R}}}
\newunicodechar{𝔽}{\ensuremath{\mathbb{F}}}
\newunicodechar{α}{\ensuremath{\alpha}}
\newunicodechar{β}{\ensuremath{\beta}}
\newunicodechar{σ}{\ensuremath{\sigma}}
\newunicodechar{τ}{\ensuremath{\tau}}
\newunicodechar{φ}{\ensuremath{\varphi}}
\newunicodechar{ψ}{\ensuremath{\psi}}
\newunicodechar{χ}{\ensuremath{\chi}}
\newunicodechar{Φ}{\ensuremath{\Phi}}
\newunicodechar{Δ}{\ensuremath{\Delta}}
\newunicodechar{₀}{\textsubscript{0}}
\newunicodechar{₁}{\textsubscript{1}}
\newunicodechar{₂}{\textsubscript{2}}
\newunicodechar{ᵏ}{\ensuremath{{}^{k}}}
\newunicodechar{ⁿ}{\ensuremath{{}^{n}}}
\newunicodechar{ʲ}{\ensuremath{{}^{j}}}
\newunicodechar{²}{\ensuremath{{}^{2}}}
\newunicodechar{³}{\ensuremath{{}^{3}}}
\newunicodechar{⁴}{\ensuremath{{}^{4}}}
\newunicodechar{⁻}{\ensuremath{{}^{-}}}
\newunicodechar{¹}{\ensuremath{{}^{1}}}

% ---- Lean listing style ----
\definecolor{leanbg}{RGB}{247,247,250}
\definecolor{leankw}{RGB}{0,90,160}
\definecolor{leancm}{RGB}{120,120,120}
\lstdefinestyle{lean}{
  backgroundcolor=\color{leanbg},
  basicstyle=\ttfamily\small,
  keywordstyle=\color{leankw}\bfseries,
  commentstyle=\color{leancm}\itshape,
  morekeywords={theorem,lemma,def,fun,Prop,Type,variable,where,
    Field,Fintype,Finite,DecidableEq,CommRing,CommSemiring,Ring,CharP,Function,
    Injective,Surjective,Bijective,by,rfl,have,rw,exact,ext,simp,intro,
    iterateFrobenius,RingHom,Commute,convert,apply,using,Nat,Coprime,gcd,
    Odd,Even,IsAPN,IsAB,walsh,Nat,obtain,induction,cases,calc,grind},
  showstringspaces=false,
  columns=fullflexible,
  extendedchars=true,
  frame=single,
  framesep=6pt,
  xleftmargin=4pt,
  aboveskip=12pt,
  belowskip=14pt,
}
\lstset{style=lean}

\newtheorem{lemma}{Lemma}
\newtheorem{theorem}{Theorem}
\newtheorem{corollary}{Corollary}

% boxed "what just happened" note
\newcommand{\note}[1]{%
  \par\smallskip
  \noindent\colorbox{leanbg}{\parbox{\dimexpr\linewidth-2\fboxsep}{\small #1}}%
  \par\smallskip}


\title{\textbf{The Doubling Shift:\\[2pt]
Cyclotomic Cosets as Orbits of a Finite-Order Multiplier Map}}
\author{}
\date{}

\begin{document}

\maketitle

\thispagestyle{empty}

\begin{abstract}
\noindent
Doubling an exponent is a dynamical system.\\
Its orbits are the cyclotomic cosets.\\
Its clock is the order of $2$ modulo $2^{\,n}-1$.

\medskip
\noindent
We strip the system down to its reusable core ---\\
multiplication by one finite-order element ---\\
and read APN-ness off as one of its invariants.

\medskip
\noindent
Everything is checked in Lean~4 / Mathlib.
\end{abstract}

\vfill

\noindent\textit{Conference paper call --- short, slow-read submission.}

\newpage


\section{The picture}

\bigskip

Fix $n$. Work modulo $m = 2^{\,n}-1$.

\bigskip

Pick one exponent $e$.

\bigskip

Double it, again and again:
\[
  e,\quad 2e,\quad 4e,\quad \dots,\quad 2^{\,n-1}e
  \pmod{2^{\,n}-1}.
\]

\bigskip

That is a sequence in $\mathbb{Z}/(2^{\,n}-1)$.

\bigskip

It comes back to $e$.

\bigskip

\note{One exponent. One periodic orbit.}

\newpage


\section{The shift}

\bigskip

Doubling once moves the whole orbit one step:
\[
  (e,\; 2e,\; 4e,\; \dots)
  \;\xmapsto{\ e\,\mapsto\, 2e\ }\;
  (2e,\; 4e,\; 8e,\; \dots).
\]

\bigskip

That is the \emph{doubling shift} $\sigma$.

\bigskip

Its $k$-th iterate is multiplication by $2^{\,k}$:
\[
  \sigma^{k}(e) = 2^{\,k}\, e .
\]

\bigskip

It is the exponent shadow of Frobenius $x \mapsto x^{2}$:
\[
  (x^{e})^{2} = x^{2e}.
\]

\bigskip

\note{Algebra of exponents on one side. A shift on the other.}

\newpage


\section{The encoding, in Lean}

\bigskip

The shift itself:

\begin{lstlisting}
def shift (m : ℕ) : ZMod m → ZMod m := fun e => 2 * e
\end{lstlisting}

\bigskip

Its $k$-th iterate is multiply-by-$2^{k}$:

\begin{lstlisting}
theorem shift_iterate (m : ℕ) (e : ZMod m) (k : ℕ) :
    (shift m)^[k] e = (2 : ZMod m) ^ k * e
\end{lstlisting}

\bigskip

\note{The iterate rule \emph{is} the doubling rule.}

\newpage


\section{Periodicity}

\bigskip

The whole engine is one identity:
\[
  2^{\,n} \equiv 1 \pmod{2^{\,n}-1}.
\]

\bigskip

So $\sigma^{\,n} = \mathrm{id}$.

\bigskip

Every exponent is periodic, and every orbit length divides $n$.

\begin{lstlisting}
theorem two_pow_eq_one (n : ℕ) :
    (2 : ZMod (2 ^ n - 1)) ^ n = 1

theorem shift_periodic (n : ℕ) (e : ZMod (2 ^ n - 1)) :
    (shift (2 ^ n - 1))^[n] e = e

theorem minimalPeriod_dvd (n : ℕ) (e : ZMod (2 ^ n - 1)) :
    Function.minimalPeriod (shift (2 ^ n - 1)) e ∣ n
\end{lstlisting}

\bigskip

\note{Orbits are the $2$-cyclotomic cosets mod $2^{\,n}-1$.}

\newpage


\section{The number-theory pearl}

\bigskip

How fast does the clock tick?

\bigskip

The order of $2$ modulo $2^{\,n}-1$ is \textbf{exactly} $n$.

\bigskip

And the same is true for \emph{any} base $b \ge 2$:
\[
  \mathrm{ord}_{\,b^{\,n}-1}(b) \;=\; n .
\]

\bigskip

The reason is one line:
\[
  b^{\,k}-1 \;<\; b^{\,n}-1
  \qquad\text{for } 0 < k < n,
\]
so $b^{\,k}\not\equiv 1$ before step $n$.

\begin{lstlisting}
theorem orderOf_base (b n : ℕ) (hb : 2 ≤ b) (hn : 1 ≤ n) :
    orderOf (b : ZMod (b ^ n - 1)) = n

theorem orderOf_two (n : ℕ) (hn : 1 ≤ n) :
    orderOf (2 : ZMod (2 ^ n - 1)) = n
\end{lstlisting}

\bigskip

\note{A crisp repunit fact: $\mathrm{ord}_{b^{n}-1}(b)=n$, base $b$ general.}

\newpage


\section{The principal orbit}

\bigskip

Start at the generator $e = 1$.

\bigskip

Its orbit is
\[
  \{\,1,\; 2,\; 4,\; \dots,\; 2^{\,n-1}\,\}.
\]

\bigskip

The clock pearl says: this orbit has \textbf{full} length $n$.

\begin{lstlisting}
theorem shift_minimalPeriod_one (n : ℕ) (hn : 1 ≤ n) :
    Function.minimalPeriod (shift (2 ^ n - 1)) 1 = n
\end{lstlisting}

\bigskip

It is the principal $2$-cyclotomic coset.

\bigskip

\note{The longest possible orbit, pinned exactly to $n$.}

\newpage


\section{The shift is invertible}

\bigskip

Since $\sigma^{\,n} = \mathrm{id}$, the inverse of $\sigma$ is just $\sigma^{\,n-1}$.

\bigskip

So $\sigma$ is a permutation of phase space:
an automorphism of $\mathbb{Z}/(2^{\,n}-1)$.

\begin{lstlisting}
theorem shift_bijective (n : ℕ) (hn : 1 ≤ n) :
    Function.Bijective (shift (2 ^ n - 1))
\end{lstlisting}

\bigskip

\note{Finite order $\Rightarrow$ invertible. An automorphism of the phase space.}

\newpage


\section{Orbits \emph{are} cyclotomic cosets}

\bigskip

Say $e$ and $d$ are on the same orbit if one reaches the other:
\[
  \mathrm{SameOrbit}\;:\quad \exists\,k,\ \sigma^{k}(d) = e .
\]

\bigskip

This is an equivalence relation (reflexive, symmetric via periodicity,
transitive).

\bigskip

Spelled out on the integers, it is the cyclotomic congruence:
\[
  e \equiv d\cdot 2^{\,j} \pmod{2^{\,n}-1}.
\]

\begin{lstlisting}
theorem sameOrbit_iff_modEq (n d e : ℕ) :
    SameOrbit (2 ^ n - 1) (d : ZMod (2 ^ n - 1))
                          (e : ZMod (2 ^ n - 1)) ↔
      ∃ j : ℕ, e ≡ d * 2 ^ j [MOD 2 ^ n - 1]
\end{lstlisting}

\bigskip

\note{Orbits partition exponent space. They are exactly the cosets.}

\newpage


\section{The reusable core}

\bigskip

None of this is really about $2$, or about fields.

\bigskip

Take any monoid and any element $c$.
Let $\sigma_c(x) = c\cdot x$.

\bigskip

Then:
\begin{itemize}
  \item $\sigma_c^{\,k}(x) = c^{\,k}\,x$;
  \item if $c^{\,d}=1$, then $\sigma_c$ is periodic of period $d$
        and a bijection with inverse $\sigma_c^{\,d-1}$;
  \item the orbit relation is an equivalence relation.
\end{itemize}

\begin{lstlisting}
theorem mulShift_iterate {M} [Monoid M] (c : M) (e : M) (k : ℕ) :
    (mulShift c)^[k] e = c ^ k * e

theorem mulShift_bijective {M} [Monoid M] {c : M} {d : ℕ}
    (hd : 1 ≤ d) (hc : c ^ d = 1) :
    Function.Bijective (mulShift c)
\end{lstlisting}

\bigskip

\note{Multiplication by a finite-order element is a finite-order permutation.
The doubling shift is just the case $c = 2$.}

\newpage


\section{The Frobenius conjugacy}

\bigskip

Why $2$? Because of Frobenius.

\bigskip

In any commutative monoid:
\[
  (x^{e})^{2} = x^{2e}.
\]

\bigskip

Doubling the exponent is post-composing with $x \mapsto x^{2}$.

\bigskip

So the power map commutes with Frobenius:
\[
  (x^{e})^{2} = (x^{2})^{e}.
\]

\begin{lstlisting}
theorem frobenius_conjugates_pow {M} [CommMonoid M]
    (e : ℕ) (x : M) : (x ^ e) ^ 2 = x ^ (2 * e)

theorem pow_comm_frobenius {M} [CommMonoid M] (e : ℕ) :
    (fun x : M => (x ^ e) ^ 2) = (fun x : M => (x ^ 2) ^ e)
\end{lstlisting}

\bigskip

\note{The doubling shift is the exponent-space shadow of Frobenius.}

\newpage


\section{APN, as a conjugacy invariant}

\bigskip

Here is the payoff slogan.

\bigskip

Any predicate on exponents that is \emph{constant on $\langle 2\rangle$-orbits}
descends to a predicate on the orbit space.

\bigskip

APN-ness of the power map $x \mapsto x^{d}$ on $\mathbb{F}_{2^{\,n}}$ is such a
predicate:
\[
  d \;\sim\; e
  \quad(\text{same orbit})
  \;\Longrightarrow\;
  \bigl(\, x^{d}\ \text{APN} \iff x^{e}\ \text{APN} \,\bigr).
\]

\bigskip

So APN-ness is a \textbf{conjugacy invariant of the doubling dynamics}:
one orbit, one verdict.

\bigskip

\note{``Same orbit, same verdict.'' The dynamical reading of cyclotomic
invariance. The field-level proof of this step lives in the project's APN
layer; here we isolate the dynamics it rests on.}

\newpage


\section{What we claim}

\bigskip

\begin{itemize}
  \item One exponent $\;\mapsto\;$ one periodic orbit under $e \mapsto 2e$.\\[6pt]
  \item $\sigma^{k} = {}$ multiply-by-$2^{k}$; \ $\sigma^{\,n}=\mathrm{id}$.\\[6pt]
  \item Orbit lengths divide $n$; orbits are the cyclotomic cosets.\\[6pt]
  \item $\mathrm{ord}_{\,b^{n}-1}(b)=n$ for every base $b\ge 2$.\\[6pt]
  \item The generator $1$ has full minimal period $n$.\\[6pt]
  \item $\sigma$ is a bijection (inverse $\sigma^{\,n-1}$).\\[6pt]
  \item The core is general: $x \mapsto c\,x$ for a finite-order $c$.\\[6pt]
  \item APN-ness is a conjugacy invariant: same orbit, same verdict.\\[6pt]
  \item All boxed statements are machine-checked, sorry-free.
\end{itemize}

\vfill

\noindent\textit{Read slowly. Doubling was a dynamical system all along.}

\end{document}
